\chapter{Consentimiento Informado del Equipo para SocialScrum}
\label{anexo:consentimiento}

\section{Consentimiento Informado para la Implementación de SocialScrum}

Yo, \underline{\hspace{6cm}} (nombre completo del miembro del equipo), en mi rol de \underline{\hspace{5cm}} (rol en el equipo: Developer, QA, etc.), declaro que he recibido una explicación clara y completa sobre la implementación del marco SocialScrum en el equipo \underline{\hspace{5cm}} (nombre del equipo).

\subsection{Comprensión del propósito de SocialScrum}
Entiendo que SocialScrum es un proceso ágil orientado a:

\begin{itemize}
    \item \textbf{Identificar, priorizar y mitigar la deuda social} del equipo, entendida como decisiones socio-técnicas que afectan nuestra productividad, calidad del software y cohesión interna.
    \item \textbf{Gestionar \gls{community-smells}} (patrones de comportamiento disfuncionales) mediante ciclos continuos de transparencia, inspección y adaptación.
    \item \textbf{Promover un entorno de trabajo sano y sostenible}, donde el bienestar del equipo y la evolución del producto avanzan de forma equilibrada.
\end{itemize}

\textbf{También entiendo que SocialScrum NO es:}

\begin{itemize}
    \item Un sistema de evaluación de desempeño individual.
    \item Un mecanismo de vigilancia o control organizacional.
    \item Terapia grupal ni un proceso de intervención psicológica.
    \item Un requisito obligatorio para mantener mi empleo. La participación es voluntaria.
\end{itemize}

\subsection{Información que será recopilada}

Entiendo que durante SocialScrum se recopilará lo siguiente:

\subsubsection{Información cuantitativa}
\begin{itemize}
    \item \textbf{Health Score}: Medición periódica (por Sprint) de los cinco valores de Scrum —Apertura, Compromiso, Respeto, Foco y Coraje— con una escala de 1 a 10.
    \item \textbf{Métricas sociales agregadas}: Patrones de comunicación, distribución del conocimiento (\gls{sna}), señales de \textit{burnout} (horas extra, cambios abruptos en la velocidad), etc.
    \item \textbf{\gls{community-smells}}: Patrones disfuncionales detectados a nivel de equipo (no individual), como \gls{cs-radio-silence}, \gls{cs-lone-wolf}, \gls{cs-black-cloud}, entre otros.
\end{itemize}

\subsubsection{Información cualitativa}

\begin{itemize}
    \item \textbf{Comentarios en eventos sociales}: Aportes realizados durante el Social Daily Standup, la Social Sprint Retrospective o las sesiones de mediación.
    \item \textbf{Entrevistas 1:1 opcionales}: Conversaciones privadas con el Social Guide para profundizar en temas relevantes.
    \item \textbf{Observaciones contextuales}: Patrones de comportamiento observados en la dinámica diaria del equipo.
\end{itemize}

\subsection{Cómo será utilizada la información}

Comprendo que la información recopilada será usada \textbf{únicamente} para:

\begin{itemize}
    \item \textbf{Mejoras internas}: Identificar \gls{community-smells}, diseñar intervenciones y medir su impacto.
    \item \textbf{Reportes agregados}: El Social Guide puede compartir únicamente datos anonimizados con el CTO o el Scrum Master (por ejemplo: ``El equipo tiene un Health Score de 5,8/10''), sin revelar identidades individuales.
    \item \textbf{Seguimiento de tendencias}: Revisar cómo evoluciona la salud social del equipo a lo largo del tiempo.
\end{itemize}

\textbf{La información NO será utilizada para:}

\begin{itemize}
    \item Evaluaciones de desempeño.
    \item Decisiones de contratación, promoción, rotación o despido.
    \item Acciones disciplinarias o llamados de atención.
    \item Determinar compensación salarial.
\end{itemize}

\subsection{Quién tendrá acceso a la información}

Entiendo que:

\begin{itemize}
    \item \textbf{El Social Guide} gestiona y protege toda la información recopilada.
    \item \textbf{El equipo completo} tendrá acceso a los reportes agregados (Health Score, \gls{community-smells}, acciones de mitigación).
    \item \textbf{El Scrum Master y el CTO/VP de Ingeniería} recibirán solo información agregada y anonimizada.
    \item \textbf{Recursos Humanos (RR. HH.)} no tendrá acceso a información individual, salvo en los casos excepcionales previstos en la sección sobre excepciones a la confidencialidad.
\end{itemize}

\subsection{Salvaguardas de confidencialidad}

Entiendo que mi privacidad está protegida mediante:

\begin{itemize}
    \item \textbf{Confidencialidad absoluta}: El Social Guide no compartirá información individual fuera del equipo.
    \item \textbf{Anonimización total}: Todos los reportes dirigidos al liderazgo serán agregados.
    \item \textbf{Separación estricta entre salud social y desempeño}: La información de SocialScrum no influirá en mi evaluación laboral.
    \item \textbf{Derecho de auditoría}: Puedo revisar cualquier reporte antes de que sea compartido fuera del equipo.
\end{itemize}

\subsection{Excepciones a la confidencialidad}

Comprendo que la confidencialidad puede suspenderse solo en dos escenarios:

\begin{itemize}
    \item \textbf{Riesgo inminente de daño}: Autolesión, daño a terceros o conducta ilegal.
    \item \textbf{Investigaciones formales por acoso o discriminación}: El Social Guide podrá compartir información relevante, idealmente con mi consentimiento, y solo aquella estrictamente relacionada con la investigación.
\end{itemize}

En estos casos, el Social Guide me informará siempre que sea legalmente posible.

\subsection{Mis derechos como participante}

Entiendo que tengo derecho a:

\begin{itemize}
    \item \textbf{Retirarme en cualquier momento}, sin consecuencias laborales.
    \item \textbf{\textit{Opt-out} parcial}: Elegir no participar en ciertos eventos sin ser penalizado(a).
    \item \textbf{Auditar reportes}: Revisar los reportes antes de ser compartidos.
    \item \textbf{Denunciar violaciones de confidencialidad} al Scrum Master, al CTO o a un auditor independiente.
    \item \textbf{Solicitar explicaciones} sobre cómo se está utilizando la información.
\end{itemize}

\subsection{Duración del consentimiento}

Este consentimiento será válido durante:

\begin{itemize}
    \item El piloto inicial de SocialScrum (mínimo 3 Sprints).
    \item Una duración indefinida, siempre y cuando:
          \begin{itemize}
              \item Las salvaguardas éticas se mantengan vigentes.
              \item Yo no haya solicitado retirarme.
              \item Al menos el 70~\% del equipo continúe apoyando SocialScrum.
          \end{itemize}
\end{itemize}

Si deseo retirarme, puedo hacerlo por escrito ante el Social Guide o el Scrum Master. A partir de ese momento, mi información dejará de ser recopilada.

\subsection{Declaración de consentimiento voluntario}

\textbf{Declaro que:}

\begin{itemize}
    \item He leído y entendido este documento.
    \item He tenido la oportunidad de hacer preguntas y todas fueron respondidas.
    \item Sé que mi participación es voluntaria y reversible.
    \item Acepto participar en SocialScrum bajo las condiciones descritas.
\end{itemize}

\vspace{1cm}

\noindent
\textbf{Firma del participante:} \underline{\hspace{6cm}} \\[0.5cm]
\textbf{Nombre completo:} \underline{\hspace{6cm}} \\[0.5cm]
\textbf{Rol en el equipo:} \underline{\hspace{6cm}} \\[0.5cm]
\textbf{Fecha:} \underline{\hspace{4cm}} \\[1cm]

\noindent
\textbf{Testigo (Social Guide):} \underline{\hspace{6cm}} \\[0.5cm]
\textbf{Nombre completo:} \underline{\hspace{6cm}} \\[0.5cm]
\textbf{Fecha:} \underline{\hspace{4cm}} \\[1cm]

\noindent
\textbf{Testigo (Scrum Master):} \underline{\hspace{6cm}} \\[0.5cm]
\textbf{Nombre completo:} \underline{\hspace{6cm}} \\[0.5cm]
\textbf{Fecha:} \underline{\hspace{4cm}} \\[1cm]

\vspace{1cm}

\subsection{Nota para el equipo}

Si menos del 70~\% de los miembros del equipo firma este consentimiento, SocialScrum NO debe implementarse. Esto asegura un apoyo mayoritario y evita que la adopción sea impuesta a quienes no están convencidos.